\documentclass{beamer}
\begin{document}

% WRAPPING-UP BONDS-----------------------------------------------------------%
\begin{frame}{Wrapping-up Bonds}
\begin{center}
    \includegraphics[scale=.3]{t}
\end{center}
\end{frame}

% A PREVIEW-------------------------------------------------------------------%
\begin{frame}{Wrapping-up Bonds}
    \begin{itemize}
        \item Some types of bond \textit{yield} simple expressions for duration 
            \begin{itemize}
                \item Zero-coupon bonds
                \item Perpetuities
            \end{itemize}
        \item Yield curve arbitrage 
        \item Practice Problems
    \end{itemize}
\end{frame}

% DURATION--------------------------------------------------------------------%
\begin{frame}{Duration}
\begin{itemize}
\item Cash flows $c_{1},c_{2},\ldots,c_{T}$ and yield $y$:
\begin{equation}
p(y)=\sum_{t=1}^{T}{\frac{c_{t}}{(1+y)^{t}}}
\end{equation}
\item The Modified Duration
    \begin{equation}
        D_{0}^{*}=\frac{1}{1+y_{0}}\cdot\frac{1}{p_{0}}\cdot\sum_{t=1}^{T}{t\cdot\frac{c_{t}}{(1+y_{0})^{t}}}
    \end{equation}
\item Two interesting bond features that affect duration
    \begin{itemize}
        \item term
        \item coupon
    \end{itemize}
\end{itemize}
\end{frame}

% PERPETUITIES----------------------------------------------------------------%
\begin{frame}{Perpetuities}
    \begin{itemize}
        \item Perpetuities: coupons and infinite term
        \item The price is
            \begin{equation}
                p(y)=\frac{c}{y}
            \end{equation}
            where $c$ is the coupon payment
        \item The modified duration is 
            \begin{equation}
                D_{0}^{*}=\frac{1}{y_{0}}
            \end{equation}
        \item The convexity is
            \begin{equation}
                C_{0}^{*}=\frac{2}{y^{2}}
            \end{equation}
    \end{itemize}
\end{frame}

% PERPETUITIES----------------------------------------------------------------%
\begin{frame}{Perpetuities: Examples}
    \begin{enumerate}
        \item (easy) A perpetuity has yield 5\%. What are its duration and convexity? 
        \item (easy) A perpetuity makes annual payments of \$1000 and has a duration of 10 years. What must it trade for?
    \end{enumerate}
\end{frame}

% ZERO-COUPON BONDS-----------------------------------------------------------%
\begin{frame}{Zero-Coupon Bonds}
    \begin{itemize}
        \item Zero-Coupons: no coupons, but variable term
        \item The price is
            \begin{equation}
                p(y)=\frac{F}{(1+y)^{T}}
            \end{equation}
            where $F$ is the face value and $T$ is the term
        \item The modified duration is 
            \begin{equation}
                D_{0}^{*}=\frac{T}{1+y_{0}}
            \end{equation}
        \item Insight: $T\:\uparrow\:\Rightarrow\:D_{0}^{*}\uparrow$
        \item The convexity is
            \begin{equation}
                C_{0}^{*}=\frac{T(T+1)}{(1+y_{0})^{2}}
            \end{equation}
    \end{itemize}
\end{frame}

% ZERO-COUPON BONDS-----------------------------------------------------------%
\begin{frame}{Zero-Coupon Bonds: Examples}
    \begin{enumerate}
        \item (easy) A ten year zero-coupon bond has yield 5\%. What are its duration and convexity? 
        \item (hard) A ten year, \$1000 zero-coupon bond has a duration of 9 years. What must it trade for?
    \end{enumerate}
\end{frame}

% EXAMPLE---------------------------------------------------------------------%
\begin{frame}{Example}
    A 30-year maturity bond has a 10\% coupon rate, paid annually. It sells today for \$800. A 20-year maturity bond has a 5\% coupon rate, also paid annually. It sells today for \$900. You forecast that in five years, 25-year maturity bonds will sell at yields to maturity of 8\% and that 15-year maturity bonds will sell at yields of 6\%. Because the yield curve is upward-sloping, the analyst believes that coupons will be invested in short-term securities at a rate of 6\%.
    \begin{itemize}
    \item What is the expected rate of return on the 30-year bond?
    \item What is the expected rate of return on the 20-year bond? 
    \end{itemize}
\end{frame}

% BOND ARBITRAGE--------------------------------------------------------------%
\begin{frame}{Example}
    The rates today on one-years and two-years are 5\% and 10\% respectively. Suppose you estimate that the yield on one-years should be 8\% next year. What arbitrage opportunity is available to you? Complete the arbitrage table. 
\end{frame}

\end{document}
